% MN submission note: keep the Springer Nature single-column sn-jnl draft style; use iicol only if the journal explicitly asks for a two-column LaTeX submission.
\documentclass[sn-nature]{sn-jnl}

\usepackage{fontspec}
\usepackage{xeCJK}
\setmainfont{TeX Gyre Termes}
\setsansfont{TeX Gyre Heros}
\setmonofont{Latin Modern Mono}
\setCJKmainfont{SimSun}
\setCJKsansfont{SimHei}
\setCJKmonofont{FangSong}

\usepackage{multirow}
\usepackage{amsmath,amssymb,amsfonts}
\usepackage{bm}
\usepackage{makecell}
\usepackage{arydshln}
\usepackage{booktabs}
\usepackage{graphicx}
\usepackage{textcomp}
\usepackage{url}
\usepackage{placeins}
\usepackage{silence}
\WarningFilter{latexfont}{Some font shapes}
\WarningFilter{latexfont}{Font shape}
\WarningFilter{latex}{Font shape}
\WarningFilter{latex}{There were undefined references}
\WarningFilter{natbib}{Citation}
\WarningFilter{natbib}{There were undefined citations}
\hbadness=10000
\vbadness=10000
\hfuzz=200pt
\vfuzz=200pt
\emergencystretch=3em

\makeatletter
\let\snorigincludegraphics\includegraphics
\RenewDocumentCommand{\includegraphics}{O{} m}{%
  \IfFileExists{#2}{%
    \snorigincludegraphics[#1]{#2}%
  }{%
    \fbox{%
      \begin{minipage}[c][0.20\textheight][c]{0.85\linewidth}%
        \centering
        Missing figure\\
        {\ttfamily\detokenize{#2}}%
      \end{minipage}%
    }%
  }%
}
\makeatother

\makeatletter
\providecommand{\headerps@out}[1]{}
\makeatother

\InputIfFileExists{values.tex}{}{%
  \typeout{Warning: values.tex not found. Manuscript-specific value macros are missing.}%
}

\newif\ifCNENEnglish
\newif\ifCNENBilingual
\newcommand{\UseCNENChinese}{\CNENEnglishfalse\CNENBilingualfalse}
\newcommand{\UseCNENEnglish}{\CNENEnglishtrue\CNENBilingualfalse}
\newcommand{\UseCNENBilingual}{\CNENEnglishfalse\CNENBilingualtrue}
\UseCNENEnglish
\newcommand{\CNEN}[2]{%
    \ifCNENBilingual
        #1%
        \if\relax\detokenize{#2}\relax\else\par #2\fi%
    \else
        \ifCNENEnglish
            #2%
        \else
            #1%
        \fi
    \fi
}

% symbols
\newcommand{\fn}{\eta} % natural frequency
\newcommand{\Sen}{S} % sensitivity on natural frequency
\newcommand{\magn}{m} % magnitude
\newcommand{\fidx}{i} % frequency index
\newcommand{\fmax}{F} % frequency index max
\newcommand{\midx}{j} % magnitude index
\newcommand{\mmax}{M} % magnitude index max
\newcommand{\pidx}{n} % sample index
\newcommand{\pmax}{N} % sample index max
\newcommand{\dataset}{{\mathcal{D}}} % dataset
\newcommand{\yseq}{{\mathbf{Y}}} % sequence of output
\newcommand{\wMAE}{\beta} % MAE loss weight
\newcommand{\lr}{\alpha} % learning rate
\newcommand{\lrdecay}{\gamma} % learning rate decay
\newcommand{\lridx}{k} % learning rate index
\newcommand{\lrperiod}{K} % learning rate period
\newcommand{\kr}{k_{\text{r}}} % rubber mold stiffness

% functions
\newcommand{\relative}[1]{\overset{\rho}{#1}} % relative value
\newcommand{\ideal}[1]{\overset{\circ}{#1}} % ideal value
\newcommand{\exam}[1]{\tilde{#1}} % measured value
\newcommand{\pred}[1]{\hat{#1}} % predicted value
\newcommand{\true}[1]{{#1}} % true value

% short names
\newcommand{\Wrecarray}{\mathcal{W}_{\text{rec}}}
\newcommand{\Winarray}{\mathcal{W}_{\text{in}}}
\newcommand{\Woutarray}{\mathcal{W}_{\text{out}}}
\newcommand{\harray}{\mathcal{H}}
\newcommand{\abase}{\exam{\Sen}  T^{2} \ideal{\fn} ^{2} \exam{\fn}  \pi^{2} \exam{\zeta}  + 2 \exam{\Sen}  T \ideal{\fn}  \exam{\fn}  \pi \ideal{\zeta}  \exam{\zeta}  + \exam{\Sen}  \exam{\fn}  \exam{\zeta}}

\raggedbottom

\begin{document}

\title[On-Board Real-Time AI Compensation System for MET]{On-Board Real-Time AI Compensation System for Nonlinear Drift in Electrochemical Seismometers}

\author[1]{\fnm{Ang} \sur{Li}}\email{liang20@mails.jlu.edu.cn}
\author*[1]{\fnm{Hongyuan} \sur{Yang}}\email{yang\_hy@jlu.edu.cn}
\author*[1]{\fnm{Huaizhu} \sur{Zhang}}\email{huaizhuzhang@jlu.edu.cn}
\author[1]{\fnm{Xintong} \sur{Dong}}
\author[1]{\fnm{Fan} \sur{Zheng}}\email{zhengfan@jlu.edu.cn}
\author[1]{\fnm{Linhang} \sur{Zhang}}\email{linxing@jlu.edu.cn}
\author[1]{\fnm{Can} \sur{Liu}}\email{canliu23@mails.jlu.edu.cn}
\author[1]{\fnm{Ruojin} \sur{Li}}\email{lirj24@mails.jlu.edu.cn}

\affil*[1]{\orgdiv{State Key Laboratory of Deep Earth Exploration and Imaging, College of Instrumentation and Electrical Engineering}, \orgname{Jilin University}, \orgaddress{\city{Changchun}, \postcode{130061}, \country{China}}}

\abstract{\CNEN{%
电化学地震检波器（MET）凭借高灵敏度、宽频带以及较好的抗倾斜和抗冲击能力，在地球物理勘探中正变得越来越重要，并被认为是下一代地震信号探测的理想方案之一。然而，由于机械、流体与电化学等多个物理过程之间存在复杂耦合，MET 会表现出显著的震级相关非线性，从而引起频率响应漂移并限制动态范围，最终影响地震数据质量和事件检测可靠性。传统力反馈方法在大震级下存在闭环振荡风险，数据驱动的深度学习方法虽然在非线性补偿中展现出高精度优势，但其密集的矩阵运算难以在资源受限的微控制器上满足实时处理需求。针对上述实时补偿与反馈稳定性问题，本文提出面向板端实时运行的前馈补偿系统 Wiener-KAN。该系统引入了 Wiener 结构“线性动态 + 静态非线性”的分解思路：在线性部分，利用目标与理想系统间的相对频率响应初始化 Wiener 线性层，实现物理先验注入和快速的频率特征提取；在非线性部分，引入带符号与对称性约束的 Kolmogorov-Arnold 网络（KAN）可训练样条激活函数，并预计算为查找表（Precomputed LUT）。此外，本文构建并开放了涵盖幅频耦合效应的三维频率响应数据集，并设计了针对频率特性漂移的幅频损失函数（AFMAE）以提高频域拟合精度。实验结果表明，在输入震级提升超过 \valInputGainExpansion~$\times$ 的条件下，Wiener-KAN 将频率漂移从 \valMainWienerFreqDrift~Hz 降低到 \valPrimaryFreqDrift~Hz，将 \valSensitivityCompareFreq~Hz 灵敏度漂移从 \valMainWienerSensDrift~$\text{V}\cdot\text{s}/\text{m}$ 降低到 \valPrimarySensDrift~$\text{V}\cdot\text{s}/\text{m}$，并将线性度误差从 \valMainWienerLinearity\% 降低到 \valPrimaryLinearity\%。实时性方面，KAN 的预计算查找表将复杂的浮点运算转化为极低延迟的内存寻址，在 STM32F405 微控制器上实现了 \valPrimaryKeilFps{} points/s 的吞吐率和 \valPrimaryKeilRamKB~KB 的极小 RAM 占用，为 MET 实时非线性补偿提供了有效方案。%
}{%
Electrochemical seismometers (MET) are gaining importance in geophysical exploration due to their high sensitivity, wide bandwidth, and strong tilt/impact resistance, and are considered a promising next-generation seismic detection solution. However, complex coupling among mechanical, fluid, and electrochemical processes induces significant magnitude-dependent nonlinearity, causing frequency response drift and limiting dynamic range, ultimately degrading seismic data quality and event detection reliability. Traditional force feedback methods risk closed-loop oscillation at large magnitudes, while data-driven deep learning methods, despite high accuracy in nonlinear compensation, struggle to meet real-time requirements on resource-constrained microcontrollers due to intensive matrix operations. To address real-time compensation and feedback-stability limitations, this paper proposes Wiener-KAN, a feedforward compensation system for on-board real-time execution. The system adopts the Wiener-model decomposition into linear dynamics and static nonlinearity: in the linear part, the Wiener linear layer is initialized with the relative frequency response between the target and ideal systems, injecting a physical prior and enabling rapid frequency feature extraction; in the nonlinear part, a sign- and symmetry-constrained Kolmogorov-Arnold network (KAN) with trainable spline activation functions is introduced and precomputed into lookup tables, replacing spline-function evaluation with low-latency table lookup during inference. Furthermore, this paper constructs and releases a three-dimensional frequency response dataset capturing amplitude-frequency coupling effects, and designs an amplitude-frequency loss (AFMAE) targeting frequency characteristic drift to improve frequency-domain fitting accuracy. Experimental results show that with input magnitude increased by over \valInputGainExpansion~$\times$, Wiener-KAN reduces frequency drift from \valMainWienerFreqDrift~Hz to \valPrimaryFreqDrift~Hz, \valSensitivityCompareFreq~Hz sensitivity drift from \valMainWienerSensDrift~$\text{V}\cdot\text{s}/\text{m}$ to \valPrimarySensDrift~$\text{V}\cdot\text{s}/\text{m}$, and linearity error from \valMainWienerLinearity\% to \valPrimaryLinearity\%. For real-time feasibility, the precomputed KAN lookup table converts complex floating-point operations into ultra-low-latency memory addressing, achieving a throughput of \valPrimaryKeilFps{} points/s and a minimal RAM footprint of \valPrimaryKeilRamKB~KB on the STM32F405 microcontroller, providing an effective solution for real-time MET nonlinear compensation.%<CNHash>9A866DB2</CNHash>
}}

\keywords{electrochemical seismometer, Kolmogorov-Arnold network, nonlinearity, frequency response, loss function, on-board real-time compensation, compensation system}

\maketitle

\section{Introduction}

\CNEN{电化学地震检波器（MET）因其高灵敏度、宽频带以及较强的抗倾斜和抗冲击性能，在地下资源勘探和地震监测中展现出良好应用前景，被认为是新一代地震传感器的重要候选\cite{10841463}。随着全球地震台网逐步实现宽频带、高动态范围的实时数据采集\cite{WOS:000853457100001}，核心传感器在线性度与响应保持能力方面面临更严格要求。然而，电化学反应与流体力学之间的复杂耦合会在 MET 的频率响应中引入显著的震级相关非线性，包括灵敏度漂移和固有频率漂移。这种漂移会扭曲振动波形检测结果，限制动态范围，并削弱地震勘探数据质量与地震事件识别可靠性。}{Electrochemical seismometers (MET) are emerging as promising candidates for next-generation seismic sensors in underground resource exploration and earthquake monitoring, owing to their high sensitivity, wide bandwidth, and strong tilt/impact resistance \cite{10841463}. As global seismic networks progressively adopt real-time data acquisition with wide bandwidth and high dynamic range \cite{WOS:000853457100001}, core sensors face stricter requirements on linearity and response fidelity. However, the complex coupling between electrochemical reactions and fluid mechanics introduces significant magnitude-dependent nonlinearity into the MET frequency response, including sensitivity drift and natural frequency drift. Such drift distorts vibration waveform detection, limits dynamic range, and degrades seismic survey data quality and event identification reliability.%<CNHash>73C9D908</CNHash>
}

\CNEN{既有研究曾尝试通过结构设计优化来提升线性度。Xu 等\cite{9481294}提出了一种采用阳极键合工艺的 MEMS 电化学地震计，通过硅-玻璃-硅三层集成电极结构提高了灵敏度和器件一致性。Sun 等\cite{sun_high-consistency_2017}开发了单硅片四微电极宽频带电化学地震计，实现了 \valXuBandwidthStart--\valXuBandwidthEnd~Hz 的工作带宽，并在 \valXuPeakSensitivityFreq~Hz 获得了 \valXuPeakSensitivity~$\text{V} \cdot \text{s}/\text{m}$ 的峰值灵敏度。该类结构优化在提升线性度方面的作用较为有限，系统通常仅在震级低于 \valXuLinearLimitMm~mm/s@\valLinearLimitRefFreq~Hz、对应加速度为 \valXuLinearLimitAcc~$\text{m}/\text{s}^2$\cite{9481294}，或低于 \valSunLinearLimitMm~mm/s@\valLinearLimitRefFreq~Hz、对应加速度为 \valSunLinearLimitAcc~$\text{m}/\text{s}^2$\cite{sun_high-consistency_2017} 时保持近似线性运行。}{Prior research has attempted to improve linearity through structural design optimization. Xu et al. \cite{9481294} proposed a MEMS electrochemical seismometer using an anodic bonding process, enhancing sensitivity and device consistency via a silicon-glass-silicon three-layer integrated electrode structure. Sun et al. \cite{sun_high-consistency_2017} developed a single-silicon-chip four-microelectrode wide-bandwidth electrochemical seismometer, achieving an operating bandwidth of \valXuBandwidthStart--\valXuBandwidthEnd~Hz and a peak sensitivity of \valXuPeakSensitivity~$\text{V} \cdot \text{s}/\text{m}$ at \valXuPeakSensitivityFreq~Hz. Such structural optimization offers limited improvement in linearity; systems typically maintain quasi-linear operation only below \valXuLinearLimitMm~mm/s@\valLinearLimitRefFreq~Hz (equivalent acceleration of \valXuLinearLimitAcc~$\text{m}/\text{s}^2$) \cite{9481294}, or below \valSunLinearLimitMm~mm/s@\valLinearLimitRefFreq~Hz (equivalent acceleration of \valSunLinearLimitAcc~$\text{m}/\text{s}^2$) \cite{sun_high-consistency_2017}.%<CNHash>BE020E37</CNHash>
}

\CNEN{力反馈是一种有效的非线性抑制手段。例如，Li 等\cite{li_electrochemical_2017}和 Sun 等\cite{WOS:000408118700046}均报道了力反馈对 MET 灵敏度漂移或输入输出非线性误差的抑制作用。然而，力反馈需要额外机械反馈单元，并受到非线性相位裕度下降和大震级闭环振荡风险限制。基于预测模型的传感器校正方法\cite{WOS:001409548400009}表明，软件前馈补偿能够在保持硬件结构简洁的同时改善传感器输出质量。}{Force feedback is an effective means of nonlinear suppression. For example, Li et al. \cite{li_electrochemical_2017} and Sun et al. \cite{WOS:000408118700046} reported reductions in MET sensitivity drift or input-output nonlinear error using force feedback. However, force feedback requires additional mechanical feedback units and is limited by reduced nonlinear phase margin and closed-loop oscillation risk at large magnitudes. Sensor calibration methods based on predictive models \cite{WOS:001409548400009} demonstrate that software-based feedforward compensation improves sensor output quality while maintaining hardware simplicity.%<CNHash>4732C298</CNHash>
}

\CNEN{这一前馈方向对 MET 具有直接意义。MET 的非线性主要表现为随震级变化的等效频率响应，而采集链路本身可以保持固定。因此，补偿器可以在不改变电化学腔体、机械拾振部件和读出电路的前提下，对已输出的时域信号进行幅值相关修正。相比依赖反馈力实时作用于机械结构，软件前馈补偿把稳定性问题转化为可测数据上的输入输出映射问题，更适合与边缘微控制器、数字采集链路和批量标定流程结合。}{This feedforward direction is directly relevant to MET sensors. The MET nonlinearity is mainly expressed as a magnitude-dependent equivalent frequency response, while the acquisition chain can remain fixed. A compensator can therefore correct the measured time-domain signal without modifying the electrochemical chamber, mechanical vibration pickup unit, or readout circuit. Compared with feedback force applied in real time to the mechanical structure, software feedforward compensation turns the stability issue into an input-output mapping problem on measured data, making it compatible with edge microcontrollers, digital acquisition chains, and batch calibration workflows.}

\CNEN{近期研究表明，当使用充分的时频域测试数据进行训练时，深度神经网络可以获得较高的传感器补偿精度\cite{HOFLER2024136528}。本文对电化学地震检波器实施覆盖 \valFreqScanStart~Hz 到 \valFreqScanEnd~Hz 全带宽的稳态测试，并将输入上限由文献\cite{WOS:000408118700046}中的 \valSunForceAcc~$\text{m}/\text{s}^2$ 扩展至 \valMagScanEnd~$\text{m}/\text{s}^2$，即提升超过 \valInputGainExpansion~$\times$。在此基础上，本文构建并开放了同时包含震级、频率以及输入输出时序的三维频率响应数据集，并显式记录幅频耦合效应。}{Recent studies show that deep neural networks can achieve high sensor compensation accuracy when trained with sufficient time- and frequency-domain test data \cite{HOFLER2024136528}. This paper performs steady-state testing on an electrochemical seismometer covering the full bandwidth from \valFreqScanStart~Hz to \valFreqScanEnd~Hz, and extends the input upper limit from \valSunForceAcc~$\text{m}/\text{s}^2$ in \cite{WOS:000408118700046} to \valMagScanEnd~$\text{m}/\text{s}^2$, an increase of over \valInputGainExpansion~$\times$. Based on this, we construct and release a three-dimensional frequency response dataset that includes magnitude, frequency, and input-output time series while explicitly recording amplitude-frequency coupling effects.%<CNHash>835747F0</CNHash>
}

\CNEN{该数据集覆盖单个频点或单个震级标定曲线之外的信息。频率维度提供低频端、共振区和高频端的响应变化，震级维度记录从小信号近似线性区到大信号非线性区的迁移过程，时域序列则保留相位、幅值和局部波形形态。三类信息共同支撑本文对频率漂移、灵敏度漂移和线性度误差的统一评价。}{The dataset is more than a calibration curve at a single frequency or magnitude. The frequency dimension records response changes at the low-frequency end, resonance region, and high-frequency end; the magnitude dimension records the transition from the small-signal quasi-linear regime to the large-signal nonlinear regime; and the time-domain sequences preserve phase, amplitude, and local waveform morphology. These three forms of information jointly support the evaluation of frequency drift, sensitivity drift, and linearity error in this work.}

\CNEN{在获得较完整测试数据之后，如何设计高效的非线性补偿算法成为核心挑战。传统信号处理辨识方法\cite{royOptimalModifiedHomotopy2023,jainModelOrderReduction2020,shawHybridAveragingHarmonic2023,helieInputOutputReduced2021}在强耦合电化学-流体系统中通常难以精确确定模型参数。数据驱动的时序模型能够提升补偿精度\cite{anandanatarajanDeepNeuralNetworkBased2023,zhaoDetectingEarlyDamages2020,thangamalarLinearisationFlowSensors2023}，但其参数量、矩阵乘加和内存开销难以满足资源受限微控制器上的实时处理要求。纯数据模型还通常忽略传感器特定的频率响应规律，而直接注入物理方程的网络\cite{wangEigenvectorBiasFourier2021}难以适配难以精确写成偏微分方程系统的 MET。}{After relatively complete test data are obtained, the core challenge becomes the design of an efficient nonlinear compensation algorithm. Conventional signal-processing identification methods \cite{royOptimalModifiedHomotopy2023,jainModelOrderReduction2020,shawHybridAveragingHarmonic2023,helieInputOutputReduced2021} usually have difficulty determining model parameters accurately in strongly coupled electrochemical-fluid systems. Data-driven time-series models can improve compensation accuracy \cite{anandanatarajanDeepNeuralNetworkBased2023,zhaoDetectingEarlyDamages2020,thangamalarLinearisationFlowSensors2023}, but their parameter counts, matrix multiply-accumulate operations, and memory costs are difficult to reconcile with real-time processing on resource-constrained microcontrollers. Purely data-driven models also usually ignore sensor-specific frequency response characteristics, while networks that directly inject physical equations \cite{wangEigenvectorBiasFourier2021} are hard to adapt to MET sensors that are difficult to formulate accurately as PDE systems.%<CNHash>790FBCE8</CNHash>
}

\CNEN{由此形成的算法目标具有三层含义。首先，补偿器需要利用 MET 在固定震级下仍近似局部线性的特征，使频率响应信息以显式形式进入网络。其次，非线性映射需要具备足够的幅值相关表达能力，以处理不同局部动态之间的残差耦合。最后，训练后的模型需要转化为低延迟、低内存占用的板端实现，使补偿过程能够跟随连续采样流逐点运行。}{The resulting algorithmic target has three implications. First, the compensator must exploit the fact that MET remains approximately locally linear at a fixed magnitude, so that frequency response information enters the network explicitly. Second, the nonlinear mapping must have sufficient magnitude-dependent expressive capacity to handle residual coupling among different local dynamics. Third, the trained model must be converted into a low-latency and low-memory on-board implementation, enabling the compensation process to run point by point along a continuous sampling stream.}

\CNEN{针对计算限制与物理先验缺失的局限，本文提出由 Wiener 线性层与 Kolmogorov-Arnold 网络相结合的板端实时 AI 补偿系统 Wiener-KAN。该系统引入传统 Wiener 模型线性动态与静态非线性分离的思路\cite{sersourNonlinearSystemIdentification2018}：在线性部分，将实测局部频率响应切片映射为由震级条件初始化的 Wiener 线性层；在非线性部分，引入带符号与对称性约束的 Kolmogorov-Arnold 网络（KAN）可训练样条激活函数\cite{schmidt-hieberKolmogorovArnoldRepresentation2021,liu2025kankolmogorovarnoldnetworks,alvesUseKolmogorovarnoldNetworks2024,livierisCKANNewApproach2024}。训练完成后，KAN 的一维样条映射被预计算为查找表，在板端推理阶段将密集浮点运算转化为低延迟内存寻址与读表操作。}{To address the limitations of computational constraints and missing physical priors, this paper proposes Wiener-KAN as an on-board real-time AI compensation system combining a Wiener linear layer with a Kolmogorov-Arnold network. The system adopts the classical Wiener model concept of separating linear dynamics from static nonlinearity \cite{sersourNonlinearSystemIdentification2018}: in the linear part, measured local frequency response slices are mapped to a Wiener linear layer initialized by magnitude conditions; in the nonlinear part, a sign- and symmetry-constrained Kolmogorov-Arnold network (KAN) with trainable spline activation functions is introduced \cite{schmidt-hieberKolmogorovArnoldRepresentation2021,liu2025kankolmogorovarnoldnetworks,alvesUseKolmogorovarnoldNetworks2024,livierisCKANNewApproach2024}. After training, the one-dimensional spline mappings in KAN are precomputed into lookup tables, transforming dense floating-point operations into low-latency memory addressing and table lookup during on-board inference.%<CNHash>F742EF80</CNHash>
}

\CNEN{据此，本文贡献围绕板端补偿系统及其数据、模型、损失和部署组件展开。三维频率响应数据集为幅频耦合提供实测证据；Wiener 线性层承载局部频率响应先验；带符号和对称性条件的 KAN 使用可训练样条函数映射剩余震级相关非线性；AFMAE 使训练目标与频率响应精度一致；预计算 LUT 将训练后的映射转化为嵌入式实现。上述组件在相同物理指标和板端协议下进行评价。}{Accordingly, the contribution of this work is organized around an on-board compensation system and its data, model, loss, and deployment components. The three-dimensional frequency response dataset supplies measured evidence for amplitude-frequency coupling; the Wiener linear layer carries the local frequency response prior; the sign- and symmetry-conditioned KAN maps the remaining magnitude-dependent nonlinearity with trainable spline functions; AFMAE aligns the training objective with frequency-response accuracy; and the precomputed LUT converts the trained mapping into an embedded implementation. These components are evaluated under the same physical metrics and on-board protocol.}



\section{Background and Principles}

\subsection{Structure and Operating Principle of MET Electrochemical Seismometers}

\begin{figure*}[ht]
    \centering
    \includegraphics[width=1.0\textwidth]{../../../ex_projects/plot/multi/fig_20_met_structure_readout_montage/data/fig_20_met_structure_readout_montage.png}
    \caption{MET sensing object, readout circuit, and large-signal nonlinear mechanism. (a) Schematic diagram of the MET structure. (b) Readout circuit of the MET. (c) Mechanism schematic of large-signal nonlinearity in MET, where mechanical geometry, fluid inertia, ion-transport coupling, and electrode kinetics jointly produce amplitude-dependent local dynamics.}
    \label{fig:met_structure_readout}
\end{figure*}

\CNEN{电化学地震检波器（MET）是一种基于分子电子转移原理检测地震振动的高灵敏度传感器。其核心组成包括电化学腔体、敏感电极、电解液以及机械拾振部件（橡胶膜与质量块），如 Fig.~\ref{fig:met_structure_readout}(a) 所示。读出电路采用 \valReadoutAmplifier{} 运算放大器将 MET 电极产生的微弱差分电流转换为电压信号，如 Fig.~\ref{fig:met_structure_readout}(b) 所示。当地震波作用于系统时，机械结构发生形变，驱动电解液流动并改变电极区域内的离子浓度分布，从而调制电极反应电流密度。}{The electrochemical seismometer (MET) is a high-sensitivity sensor that detects seismic vibrations based on molecular electron transfer. Its core components include an electrochemical chamber, sensing electrodes, electrolyte, and a mechanical vibration pickup unit (rubber membrane and mass body), as shown in Fig.~\ref{fig:met_structure_readout}(a). The readout circuit uses an \valReadoutAmplifier{} operational amplifier to convert the weak differential current generated by the MET electrodes into a voltage signal, as shown in Fig.~\ref{fig:met_structure_readout}(b). When a seismic wave acts on the system, the mechanical structure deforms, driving electrolyte flow and altering the ion concentration distribution in the electrode region, thereby modulating the electrode reaction current density.%<CNHash>6BB89E74</CNHash>
}

\CNEN{在小信号条件下，MET 的频率响应可由机械拾振子系统和电化学转换子系统串联得到，相关机械与电化学小信号子模型已在既有 MET 研究中给出\cite{WOS:000318036400036,agafonov2015,sun_influence_2011,sun_numerical_2012,kozlov_frequency_2002}。机械子系统把外部振动转换为流道内电解液体积流量，电化学子系统再把流量诱导的离子浓度扰动转换为电极电流。由此可用机械二阶高通环节与电化学一阶低通环节共同描述理想小信号响应，其中机械参数控制低频端和共振区，电极附近离子扩散控制高频端。完整推导见 Supplementary Note 1。}{Under small-signal conditions, the MET frequency response can be obtained by cascading the mechanical vibration pickup subsystem and the electrochemical conversion subsystem, and the corresponding mechanical and electrochemical small-signal submodels have been established in prior MET studies \cite{WOS:000318036400036,agafonov2015,sun_influence_2011,sun_numerical_2012,kozlov_frequency_2002}. The mechanical subsystem converts external vibration into volumetric flow rate of the electrolyte in the flow channel, and the electrochemical subsystem converts the flow-induced ion concentration perturbation into electrode current. The ideal small-signal response can therefore be described by a mechanical second-order high-pass stage and an electrochemical first-order low-pass stage, where the mechanical parameters govern the low-frequency end and resonance region, while ion diffusion near the electrodes controls the high-frequency end. The full derivation is given in Supplementary Note 1.%<CNHash>F8EDC391</CNHash>
}

\CNEN{这一小信号模型为本文指标和模型设计提供了共同参照。固有频率可以概括低频端与共振区的动态迁移，灵敏度可以描述通带增益变化，阻尼比则反映响应曲线宽度和峰值形态。虽然大信号状态下完整多物理过程难以写成一个精确闭式模型，但这些等效参数能够把复杂耦合响应压缩为可测、可比较的频率响应特征。}{This small-signal model provides a common reference for the metrics and model design used later in this paper. The natural frequency summarizes the dynamic shift at the low-frequency end and resonance region, the sensitivity describes passband gain variation, and the damping ratio reflects the response-curve width and peak shape. Although the complete multiphysics process under large-signal operation is difficult to express as an exact closed-form model, these equivalent parameters compress the coupled response into measurable and comparable frequency-response features.}

\CNEN{当输入震级升高时，小信号线性化条件逐渐偏离实际工作状态。机械几何非线性、流体惯性非线性、离子输运乘法耦合和电极反应动力学指数非线性会共同改变局部频率响应，并表现为等效固有频率、灵敏度和阻尼比随输入震级迁移，如 Fig.~\ref{fig:met_structure_readout}(c) 所示。该幅频耦合机理的公式展开见 Supplementary Note 1。}{As input magnitude increases, the small-signal linearization condition gradually deviates from the actual operating state. Mechanical geometric nonlinearity, fluid inertial nonlinearity, multiplicative ion-transport coupling, and the exponential nonlinearity of electrode reaction kinetics jointly alter the local frequency response, appearing as magnitude-dependent shifts in equivalent natural frequency, sensitivity, and damping ratio, as shown in Fig.~\ref{fig:met_structure_readout}(c). The formula expansion of this amplitude-frequency coupling mechanism is provided in Supplementary Note 1.%<CNHash>59CE5B77</CNHash>
}

\CNEN{因此，本文把固定震级下的局部频率响应视为连接物理机理与学习模型的中间表示。该表示避免直接求解完整电化学-流体非线性方程，同时保留了 MET 在不同输入震级下表现出的主要动态变化。Wiener 线性层利用这一局部线性特征，把多个震级条件下的相对传递函数组织成并行线性核。}{This paper therefore treats the local frequency response at a fixed magnitude as an intermediate representation linking the physical mechanism and the learning model. This representation avoids directly solving the complete electrochemical-fluid nonlinear equations while preserving the dominant dynamic changes observed under different input magnitudes. The subsequent Wiener linear layer uses this local linear characteristic to organize relative transfer functions under multiple magnitude conditions into parallel linear kernels.}



\section{Experimental Design and Dataset}

\CNEN{精确评估并补偿 MET 的幅值相关非线性漂移，要求对其动态响应进行全面表征。传统单维度频率扫描\cite{WOS:000492361300009}无法揭示由输入幅值引起的系统动态变化，因此本文在频率与震级联合空间中进行二维扫描，并保留完整输入输出时序，以同时记录波形畸变、相位偏移与局部动态细节。}{Precise evaluation and compensation of the amplitude-dependent nonlinear drift of MET sensors require comprehensive characterization of their dynamic response. Conventional single-dimensional frequency sweeps \cite{WOS:000492361300009} cannot reveal system dynamics induced by input magnitude variations. This paper therefore performs two-dimensional sweeps in the joint frequency--magnitude space and preserves complete input-output time series to record waveform distortion, phase shift, and local dynamic details simultaneously.%<CNHash>E53805D4</CNHash>
}

\CNEN{该实验设计把非线性表征、训练样本构建和结果评价放在同一频率--震级矩阵内完成。每一个工况既提供未补偿 MET 的实测响应，也提供与相同输入条件对应的理想线性参考输出。因此，模型训练阶段学习的是从实测输出到目标线性输出的补偿映射，结果评价阶段则重新回到同一物理矩阵中检查频率响应漂移是否被抑制。}{This experimental design places nonlinear characterization, training-sample construction, and result evaluation within the same frequency--magnitude matrix. Each operating condition provides both the measured uncompensated MET response and the ideal linear reference output under the same input condition. Thus, the training stage learns a compensation mapping from the measured output to the target linear output, while the evaluation stage returns to the same physical matrix to examine whether frequency response drift is suppressed.}

\subsection{Experimental Setup}

\begin{figure*}[ht]
    \centering
    \includegraphics[width=1.0\textwidth]{../../../ex_projects/plot/multi/fig_21_experimental_setup_dataset_workflow_montage/data/fig_21_experimental_setup_dataset_workflow_montage.png}
    \caption{Experimental setup, dataset construction workflow, and nonlinear response characterization. (a) Photograph of the experimental setup. (b) Illustrated dataset construction and preprocessing workflow. Controlled vibration-table sweeps define the frequency--magnitude operating grid, measured MET outputs are paired with the ideal linear reference, each clipped steady-state sequence is divided into temporal training/validation halves, and the input/target sequences are min-max normalized to [-1, 1] before tensor packaging. (c) Measured nonlinear frequency responses of the uncompensated MET under the frequency--magnitude sweep.}
    \label{fig:experimental_setup_dataset_workflow}
\end{figure*}

\CNEN{如 Fig.~\ref{fig:experimental_setup_dataset_workflow}(a) 所示，振动台按照频率--震级矩阵产生正弦稳态振动信号，MET 将其转换为电信号，再由数据采集系统记录。被测样品为未采用力反馈的 \valTestSampleModel{}\cite{RSensors2024} 电化学腔体。测试矩阵覆盖 \valFreqScanPoints{} 个频率序列与 \valMagScanPoints{} 个震级序列，频率范围为 \valFreqScanStart--\valFreqScanEnd~Hz，震级范围为 \valMagScanStart--\valMagScanEnd~$\text{m}/\text{s}^2$，采样频率为 \valSamplingFreq~Hz。所有频率--震级工况采用相同的激励、采集和参考通道设置，数据记录阶段仅改变目标频率与震级。详细实验参数见 Supplementary Note 3。}{As shown in Fig.~\ref{fig:experimental_setup_dataset_workflow}(a), the vibration table generates steady-state sinusoidal signals according to the frequency--magnitude matrix. The MET sensor converts these signals into electrical outputs, which are recorded by the data acquisition system. The device under test is a \valTestSampleModel{} electrochemical chamber without force feedback \cite{RSensors2024}. The test matrix covers \valFreqScanPoints{} frequency sequences and \valMagScanPoints{} magnitude sequences, with a frequency range of \valFreqScanStart--\valFreqScanEnd~Hz, a magnitude range of \valMagScanStart--\valMagScanEnd~$\text{m}/\text{s}^2$, and a sampling frequency of \valSamplingFreq~Hz. All frequency--magnitude conditions share identical excitation, acquisition, and reference channel settings; only the target frequency and magnitude are changed during recording. Detailed experimental parameters are listed in Supplementary Note 3.%<CNHash>579DE307</CNHash>
}

\CNEN{在采集过程中保持激励、采集和参考通道设置一致，可以减少由测试链路变化引入的额外差异。这样，频率响应曲线随震级产生的变化主要对应 MET 本体的幅值相关动态，外部采集条件在各工况间保持一致。该设置也使低幅值参考系统、未补偿响应和补偿后输出能够在同一坐标下比较。}{Keeping the excitation, acquisition, and reference channel settings consistent during acquisition reduces extra differences introduced by changes in the test chain. The variation of the frequency response curves with magnitude is therefore mainly attributed to the amplitude-dependent dynamics of the MET itself, while the external acquisition conditions remain consistent across operating points. This setting also allows the low-amplitude reference system, uncompensated response, and compensated output to be compared in the same coordinates.}

\subsection{Dataset}

\CNEN{由上述实验得到的数据集包含 \valDatasetOperatingPoints{} 个频率--震级工况，每个工况保存 \valDatasetPointsPerRecord{} 个稳态采样点。数据集以成对时域序列作为基本样本，既保留幅频特征点，也保留端到端补偿器训练所需的逐点波形、相位和局部动态信息。数据集可以表示为一个 $\fmax \times \mmax$ 的矩阵 $\dataset$：}{The dataset obtained from the above experiments contains \valDatasetOperatingPoints{} frequency--magnitude operating conditions, with \valDatasetPointsPerRecord{} steady-state sampling points saved per condition. It uses paired time-domain sequences as basic samples, preserving both amplitude-frequency feature points and the pointwise waveform, phase, and local dynamic information required for end-to-end compensator training. The dataset can be represented as an $\fmax \times \mmax$ matrix $\dataset$:%<CNHash>54B3DD38</CNHash>
}

\begin{equation}
    \dataset = \left\{ \left(\exam{\yseq}_{\fidx,\midx}, \ideal{\yseq}_{\fidx,\midx} \right) \ \middle| \ \fidx = 1, \dots, \fmax, \midx = 1, \dots, \mmax \right\}.
    \label{eq:dataset}
\end{equation}

\CNEN{其中，$\exam{\yseq}_{\fidx,\midx}$ 为实测 MET 原始输出序列，$\ideal{\yseq}_{\fidx,\midx}$ 为相同频率和震级条件下理想线性系统的输出序列。目标序列由同一激励条件下的理想二阶线性参考系统生成，该参考系统来自低幅值输入时的系统。数据预处理流程包括稳态片段截取、零均值校正、频率与震级条件编码、每个工况内的训练/验证时域划分、输入/目标序列归一化，以及基于相同划分的指标重算，如 Fig.~\ref{fig:experimental_setup_dataset_workflow}(b) 所示。详细协议见 Supplementary Note 3。}{Here, $\exam{\yseq}_{\fidx,\midx}$ is the measured raw MET output sequence, and $\ideal{\yseq}_{\fidx,\midx}$ is the output sequence of an ideal linear system under the same frequency and magnitude conditions. The target sequence is generated by an ideal second-order linear reference system identified from the low-amplitude input regime. The preprocessing pipeline includes steady-state segment extraction, zero-mean correction, frequency and magnitude condition encoding, per-condition temporal training/validation splitting, input/target sequence normalization, and metric recomputation based on the same split, as shown in Fig.~\ref{fig:experimental_setup_dataset_workflow}(b). The detailed protocol is provided in Supplementary Note 3.%<CNHash>86ADD752</CNHash>
}

\CNEN{成对序列的定义使时域误差和频域误差可以从同一批样本中计算。逐点波形误差用于检查补偿输出是否贴合理想参考序列，频率响应重算则用于检查同一模型在不同震级下是否保持稳定的固有频率和灵敏度。训练/验证划分在每个频率--震级工况内执行：截取后的稳态序列被分为两个时域半段，并分别用于训练和验证。这样，每个工况都参与模型训练和评价，评价结果仍覆盖完整的频率--震级矩阵。}{The paired-sequence definition allows time-domain and frequency-domain errors to be computed from the same samples. Pointwise waveform error checks whether the compensated output follows the ideal reference sequence, while frequency-response recomputation examines whether the same model maintains stable natural frequency and sensitivity across magnitudes. The training/validation split is performed within each frequency--magnitude operating condition: the clipped steady-state sequence is divided into two temporal halves that are assigned to training and validation. Thus, every operating condition participates in both model training and evaluation, and the evaluation still covers the complete frequency--magnitude matrix.}

\subsection{Nonlinear Frequency Response Characteristics of MET}

\CNEN{Fig.~\ref{fig:experimental_setup_dataset_workflow}(c) 展示了未补偿 MET 在不同输入震级下的实测频率响应。在 \valMagScanStart--\valMagScanEnd~$\text{m}/\text{s}^2$ 震级扫描范围内，固有频率覆盖 \valMetFnMin--\valMetFnMax~Hz，\valSensitivityCompareFreq~Hz 灵敏度覆盖 \valMetSensMin--\valMetSensMax~$\text{V}\cdot\text{s}/\text{m}$。这反映了系统由于电化学与流体力学耦合而产生的显著幅频耦合漂移现象。}{Fig.~\ref{fig:experimental_setup_dataset_workflow}(c) shows the measured frequency response of the uncompensated MET under different input magnitudes. Across the \valMagScanStart--\valMagScanEnd~$\text{m}/\text{s}^2$ magnitude sweep, the fitted natural frequency spans \valMetFnMin--\valMetFnMax~Hz, and the \valSensitivityCompareFreq~Hz sensitivity spans \valMetSensMin--\valMetSensMax~$\text{V}\cdot\text{s}/\text{m}$. This reflects the significant amplitude-frequency coupling drift caused by the electrochemical-fluid system.%<CNHash>7C3DE96D</CNHash>
}

\CNEN{尽管 MET 的全量程频率响应表现出显著的震级依赖性，但在固定单一震级下，其实测频率响应仍可近似为局部线性二阶系统。基于这一局部线性特征，三支路并联 Wiener 等效模型能够通过组合多个线性动态核与静态非线性映射来表征幅值升高时的固有频率右移与 \valWienerParallelGainFreq~Hz 灵敏度抬升；其中心频率 MAE 为 \valWienerParallelCfMae~Hz，\valWienerParallelGainFreq~Hz 灵敏度 MAE 为 \valWienerParallelGainMae。详细等效建模过程见 Supplementary Note 2。}{Although the full-range frequency response of the MET exhibits significant magnitude dependence, its measured frequency response at a fixed magnitude can still be approximated as a local linear second-order system. Based on this local linear characteristic, a three-branch parallel Wiener equivalent model can combine multiple linear dynamic kernels with static nonlinear mappings to represent the rightward shift of natural frequency and the sensitivity increase at \valWienerParallelGainFreq~Hz as the input magnitude rises; it yields a center-frequency MAE of \valWienerParallelCfMae~Hz and a sensitivity MAE of \valWienerParallelGainMae at \valWienerParallelGainFreq~Hz. The detailed equivalent modeling process is provided in Supplementary Note 2.%<CNHash>8C5276BD</CNHash>
}



\section{Proposed Method}

\begin{figure*}[ht]
    \centering
    \includegraphics[width=1.0\textwidth]{../../../ex_projects/plot/multi/fig_23_wiener_kan_framework_local_transfer_montage/data/fig_23_wiener_kan_framework_local_transfer_montage.png}
    \caption{Wiener-KAN compensation framework and magnitude-conditioned local transfer prior. (a) Schematic of the Wiener-KAN compensation framework. (b) Magnitude-conditioned local frequency responses identified from fixed-magnitude MET measurements. (c) Corresponding local relative transfer functions used to initialize the Wiener linear layer.}
    \label{fig:wiener_kan_framework}
    \label{fig:FRIRNN_3D_response_slices}
\end{figure*}

\CNEN{本文提出的板端实时 AI 补偿系统如 Fig.~\ref{fig:wiener_kan_framework}(a) 所示。Wiener-KAN 将补偿问题表述为神经 Wiener 系统：用多核 Wiener 线性层同时提取各工作点下的频响特征，以保留局部物理先验；再采用多层、多输入 KAN 网络，以端到端学习方式拟合复杂动态余项与跨支路串扰，从而实现宽频带、全量程的整体映射重构。相比传统力反馈方案，该前馈结构不需要额外机械反馈装置，因此能够避免力反馈系统中由非线性引起的闭环振荡问题。并联 Wiener 逆模型的解析困难见 Supplementary Note 4。}{The proposed on-board real-time AI compensation system is illustrated in Fig.~\ref{fig:wiener_kan_framework}(a). Wiener-KAN casts the compensation problem into a neural Wiener system: a multi-kernel Wiener linear layer extracts frequency response features at different operating points to retain local physical priors, and a multi-layer, multi-input KAN network then learns the complex dynamic residuals and cross-branch crosstalk end to end, enabling global mapping reconstruction over a wide bandwidth and full dynamic range. Compared with traditional force feedback approaches, this feedforward structure requires no additional mechanical feedback unit, thereby avoiding closed-loop oscillation caused by nonlinearity in force feedback systems. The analytical difficulty of the parallel Wiener inverse model is given in Supplementary Note 4.%<CNHash>DFEEC505</CNHash>
}

\CNEN{这种结构选择来自前述局部线性观察。若只使用单一全局滤波器，模型难以同时覆盖小震级下的参考响应和大震级下的峰频迁移；若只使用黑箱时序网络，频率响应先验无法以稳定形式进入模型。Wiener-KAN 将局部线性动态与可训练非线性映射串联，使频域先验先对输入序列进行结构化投影，再由 KAN 处理剩余的幅值相关映射。}{This structural choice follows the local-linearity observation described above. A single global filter has difficulty covering both the reference response at small magnitudes and the peak-frequency shift at large magnitudes, while a black-box time-series network lacks a stable pathway for injecting frequency response priors. Wiener-KAN cascades local linear dynamics with a trainable nonlinear mapping, allowing the frequency-domain prior to first project the input sequence into structured features and then allowing KAN to process the remaining magnitude-dependent mapping.}

\CNEN{在实际推理链路中，输入序列首先经过多核 Wiener 线性层，得到与不同震级局部动态对应的一组特征序列。随后，KAN 在这些特征上执行多输入非线性组合，把局部频率响应差异、幅值相关标定曲线以及跨支路残差统一映射到目标线性输出。该流程保持了 Wiener 模型的物理可解释结构，同时避免手工推导全量程解析逆模型。}{In the actual inference path, the input sequence first passes through the multi-kernel Wiener linear layer, producing a group of feature sequences associated with local dynamics at different magnitudes. KAN then performs multi-input nonlinear combination on these features, mapping local frequency response differences, magnitude-dependent calibration curves, and cross-branch residuals to the target linear output. This workflow preserves the physically interpretable structure of the Wiener model while avoiding manual derivation of a full-range analytical inverse model.}

\subsection{KAN-Based Nonlinear Mapping with Sign and Symmetry Constraints}

\CNEN{为了构建兼具结构化物理约束和非线性建模能力的补偿器，本文通过 KAN 引入可训练 B 样条激活函数，利用其分段局部逼近特性直接学习传感器标定曲线。与 MLP 使用固定形状激活函数进行间接拟合不同，本文采用的带符号与对称性约束的 KAN 非线性映射通过可学习 B 样条系数来优化激活函数形状，从而提升单神经元对传感器非线性的建模能力\cite{livierisCKANNewApproach2024}。}{To construct a compensator that integrates structured physical constraints with nonlinear modeling capability, this paper introduces trainable b-spline activation functions via KAN, using their piecewise local approximation property to directly learn the sensor calibration curve. Unlike MLPs that use fixed-shape activation functions for indirect fitting, the sign- and symmetry-constrained KAN nonlinear mapping adopted here optimizes the activation function shape through learnable b-spline coefficients, thereby enhancing the single-neuron modeling capacity for sensor nonlinearity~\cite{livierisCKANNewApproach2024}.%<CNHash>9012E705</CNHash>
}

\begin{equation}
    \phi(x) = \text{sign}(x) \sum_{k=1}^{K} c_k B_k(|x|), \quad B_i(x) \geq 0 \ \forall x \in \mathbb{R}, \quad c_k \geq 0.
    \label{fun:KAN_activation}
\end{equation}

\CNEN{其中，$B_k(x)$ 表示 B 样条基函数，$c_k$ 表示可训练权重。奇对称性 $\phi(-x)=-\phi(x)$ 对应 MET 对正负振动输入的方向一致响应，正值性 $\phi(x)>0\ (x>0)$ 对应正输入区间内的单调方向保持。单神经元示例和多层 KAN 的完整表达见 Supplementary Note 4。}{Here, $B_k(x)$ denotes the b-spline basis function and $c_k$ represents the trainable weights. Odd symmetry $\phi(-x)=-\phi(x)$ corresponds to the direction-consistent response of the MET sensor to positive and negative vibration inputs, while positivity $\phi(x)>0\ (x>0)$ corresponds to monotonic direction preservation within the positive input interval. The single-neuron example and the full expression of multi-layer KAN are provided in Supplementary Note 4.%<CNHash>581D10FE</CNHash>
}

\CNEN{这些条件使 KAN 的函数空间与 MET 输出的方向性保持一致。对于正向和反向振动输入，传感器读出极性随输入方向变化；对于同一方向内逐渐增大的幅值，补偿映射需要维持稳定的方向关系。把奇对称性和正值性写入激活函数形式，可以减少在数据稀疏幅值区间出现符号翻转或局部不合理弯折的可能性。}{These conditions align the function space of KAN with the directional behavior of the MET output. For positive and negative vibration inputs, the sensor readout polarity changes with the input direction; within one direction, the compensation mapping should maintain a stable directional relation as the magnitude increases. Embedding odd symmetry and positivity into the activation-function form reduces the possibility of sign reversal or locally implausible bending in magnitude intervals with sparse data.}

\subsection{Magnitude-Conditioned Local Transfer Functions for Wiener Layer Initialization}

\CNEN{本文将 Wiener 结构中的线性动态部分实现为一个可训练的 Wiener 线性层，并利用震级条件局部传递函数完成初始化。在局部工作条件下，MET 可以被近似为线性系统。因此，对于每一个固定震级，本文从测量数据中辨识其对应的局部线性频率响应，并依据线性系统相对传递函数方法计算补偿器在相同工作点下对应的局部相对传递函数。通过将这些随震级变化的局部传递函数作为 Wiener 线性层的初始化依据，本文把信号处理领域知识引入到模型内部，以增强其频率相关补偿能力。}{This paper implements the linear dynamic part of the Wiener structure as a trainable Wiener linear layer, initialized using magnitude-conditioned local transfer functions. Under local operating conditions, the MET can be approximated as a linear system. Therefore, for each fixed magnitude, this paper identifies its corresponding local linear frequency response from measurement data and calculates the corresponding local relative transfer function of the compensator at the same operating point according to the relative transfer function method for linear systems. By using these magnitude-varying local transfer functions as the initialization basis for the Wiener linear layer, this paper introduces signal-processing domain knowledge into the model to enhance its frequency-dependent compensation capability.%<CNHash>DE232EA5</CNHash>
}

\begin{equation}
    \text{Wiener-KAN}(x[n]) = \operatorname{KAN}_{\text{con}} \circ \mathrm{WienerLayer}(x[n]).
    \label{eq:wiener_kan}
\end{equation}

\CNEN{针对不同震级工作点，本文基于提取出的差分方程参数构造 \mmax~个相互独立的离散线性核，并令这些局部线性滤波核并行作用于同一输入信号，形成单输入多输出（SIMO）的多核 Wiener 线性特征提取层。训练过程中，该线性层的权重被冻结，从而保留传感器在不同震级下的客观物理传递先验。局部传递函数切片见 Fig.~\ref{fig:wiener_kan_framework}(b,c)，离散差分映射见 Supplementary Note 4。}{For different magnitude operating points, this paper constructs \mmax~mutually independent discrete linear kernels from the extracted difference equation parameters and applies these local linear filtering kernels in parallel to the same input signal, forming a single-input multiple-output (SIMO) multi-kernel Wiener linear feature extraction layer. During training, the weights of this linear layer are frozen, preserving the objective physical transfer prior of the sensor at different magnitudes. The local transfer function slices are shown in Fig.~\ref{fig:wiener_kan_framework}(b,c), and the discrete difference mapping is provided in Supplementary Note 4.%<CNHash>594D498F</CNHash>
}

\CNEN{冻结该线性层具有明确作用。局部传递函数来自实测频率响应切片，代表在各震级附近已经辨识出的客观动态；训练过程中保留这些线性核，可以避免网络为了降低逐点损失而任意改变频率响应骨架。KAN 后级则保留可训练性，用于吸收局部线性近似之外的动态余项、跨支路耦合以及传感器标定曲线中的非线性部分。}{Freezing this linear layer has a specific role. The local transfer functions are derived from measured frequency response slices and represent the objective dynamics identified near each magnitude. Preserving these linear kernels during training prevents the network from arbitrarily changing the frequency response skeleton merely to reduce pointwise loss. The KAN back-end remains trainable to absorb dynamic residuals beyond the local linear approximation, cross-branch coupling, and the nonlinear part of the sensor calibration curve.}

\subsection{AFMAE Loss for Amplitude-Frequency Response}

\CNEN{传统平均绝对误差（MAE）和均方误差（MSE）主要衡量逐点误差累积，并且对高频随机噪声往往过于敏感。对于非线性系统建模而言，研究者更关注信号整体频率响应特性的保持。因此，本文设计幅频平均绝对误差（AFMAE），用于在频域量化补偿输出与理想输出之间的差异。AFMAE 基于同一频率和震级条件下的输出能量统计计算对数幅值差异，其形式为}{Conventional mean absolute error (MAE) and mean squared error (MSE) primarily measure pointwise error accumulation and tend to be overly sensitive to high-frequency random noise. For nonlinear system modeling, preserving the overall frequency response characteristics of the signal is more important. This paper therefore designs the amplitude-frequency mean absolute error (AFMAE) to quantify the difference between the compensated output and the ideal output in the frequency domain. AFMAE computes the logarithmic amplitude discrepancy from the output energy statistic under the same frequency and magnitude condition as%<CNHash>D08FA343</CNHash>
}

\begin{equation}
    \begin{aligned}
        \mathcal{L}_{\text{AFMAE}} & =  \frac{1}{\fmax\mmax} \sum_{\fidx=1}^{\fmax} \sum_{\midx=1}^{\mmax}        \left|
        \log \sum_{\pidx=0}^{\pmax-1}y_{\fidx,\midx}[n]^2
        - \log \sum_{\pidx=0}^{\pmax-1}\hat{y}_{\fidx,\midx}[n]^2
        \right|.
    \end{aligned}
    \label{eq:AFMAE_loss2}
\end{equation}

\CNEN{其中，$\fmax$ 为参与计算的频率样本数，$\mmax$ 为参与计算的震级样本数。该方法只需进行一次能量统计，便可较高效地刻画系统的幅频特性。在实际训练中，MAE 与 AFMAE 的加权组合能够得到较优结果，本文在总损失函数中引入权重超参数 $\wMAE$：}{Here, $\fmax$ is the number of frequency samples, and $\mmax$ is the number of magnitude samples. This method characterizes the amplitude-frequency response efficiently with a single energy statistic. In practice, a weighted combination of MAE and AFMAE yields superior results, and this paper introduces a weight hyperparameter $\wMAE$ into the total loss function:%<CNHash>A1749DDF</CNHash>
}

\begin{equation}
    \mathcal{L}_{\text{total}}
    = (1-\wMAE)\,\mathcal{L}_{\text{MAE}}
    + \wMAE\,\mathcal{L}_{\text{AFMAE}}.
    \label{eq:TotalLoss}
\end{equation}

\CNEN{AFMAE 的能量-幅值推导和对数形式化简见 Supplementary Note 5。}{The energy-amplitude derivation and logarithmic-form simplification of AFMAE are given in Supplementary Note 5.%<CNHash>473DCF69</CNHash>
}

\CNEN{联合损失的设计对应两类误差来源。MAE 直接限制时间序列逐点偏差，有助于保持相位和局部波形一致；AFMAE 从同一工况下的能量统计出发，使训练过程关注补偿后幅频响应与目标响应之间的差异。两者共同使用时，模型既需要输出正确的时域波形，也需要在频率--震级矩阵上保持稳定的幅值响应。}{The joint loss corresponds to two sources of error. MAE directly limits pointwise time-series deviation, helping preserve phase and local waveform consistency. AFMAE starts from the energy statistic under the same operating condition, making training sensitive to the discrepancy between the compensated amplitude-frequency response and the target response. When used together, the model must produce the correct time-domain waveform and maintain a stable amplitude response over the frequency--magnitude matrix.}

\CNEN{这一损失设计也与本文的评价指标保持一致。频率漂移和灵敏度漂移来自补偿后频率响应的重算，线性度误差来自带内输入--输出幅值关系；如果训练目标只关注逐点误差，模型可能在局部波形上获得较小偏差，却仍保留幅值相关频响差异。AFMAE 将这种差异提前纳入优化过程，从而使训练目标更接近最终物理评价。}{This loss design is also consistent with the evaluation metrics used in this paper. Frequency drift and sensitivity drift are obtained by recomputing the compensated frequency response, while linearity error is derived from the in-band input-output amplitude relation. If the training objective focuses only on pointwise error, a model may obtain small local waveform deviation while retaining magnitude-dependent frequency-response differences. AFMAE brings this discrepancy into the optimization process, making the training objective closer to the final physical evaluation.}

\subsection{Precomputed LUT Inference for Real-Time Embedded Deployment}

\CNEN{在典型地震监测应用中，连续高频采样要求补偿算法的单步推理延迟小于采样周期。KAN 的在线推理代价主要来自每条边上的可学习激活函数求值。为满足实时要求，本文在训练完成后将该激活函数离线采样为预计算查找表（Precomputed LUT），使在线推理阶段把连续样条求值转换为低延迟内存地址映射、读表以及可选一阶线性插值。}{In typical seismic monitoring applications, continuous high-frequency sampling requires the single-step inference latency of the compensation algorithm to be smaller than the sampling period. The online inference cost of KAN mainly arises from evaluating the learnable activation function on each edge. To meet real-time requirements, this paper offline-samples this activation function into a precomputed lookup table after training, converting continuous spline evaluation into low-latency memory address mapping, table lookup, and optional first-order linear interpolation during online inference.%<CNHash>F1A38F45</CNHash>
}

\CNEN{该部署路径与 MLP 的计算结构存在本质差异。MLP 层中的权重通过矩阵乘加参与每个采样点的在线计算，而 KAN 中的边权重由可训练函数表达；完成训练后，函数形状被预计算进表值，在线路径主要围绕表地址映射、读表和可选插值展开。由此，B 样条网格数和阶数主要影响离线函数表达能力与表值生成，导出实现的实时性通过目标板吞吐率直接评价。LUT 网格、最近邻查询和线性插值公式见 Supplementary Note 6。}{This deployment path differs fundamentally from the computational structure of an MLP. In an MLP layer, weights participate in online computation for every sample through matrix multiplication-addition, while edge weights in KAN are expressed by trainable functions. After training, the function shapes are precomputed into table values, and the online path mainly revolves around table address mapping, lookup, and optional interpolation. Consequently, the B-spline grid size and order primarily affect offline function expressiveness and table value generation, while the real-time feasibility of the exported implementation is directly evaluated by target-board throughput. The LUT grid, nearest-neighbor query, and linear-interpolation formulas are given in Supplementary Note 6.%<CNHash>CDBBC919</CNHash>
}
\FloatBarrier

\CNEN{嵌入式实现中，LUT 表与归一化常数、Wiener 线性核和输出缩放参数一起写入 C 端只读数据区。在线阶段按照固定采样窗口逐点执行线性递归和 KAN 查询路径，因此实际速度同时受到表访问、函数调用和编译器优化影响。本文使用 QEMU-MAE、KEIL-MAE、KEIL speed、Flash 和 RAM 同时评价导出实现，使数值一致性和实时性由同一 C 内核验证。}{In the embedded implementation, the LUT tables are written into the read-only C data region together with normalization constants, Wiener linear kernels, and output scaling parameters. During online operation, the linear recursion and KAN query path are executed point by point within a fixed sampling window, so the actual speed is affected by table access, function calls, and compiler optimization. This paper evaluates the exported implementation using QEMU-MAE, KEIL-MAE, KEIL speed, Flash, and RAM, allowing numerical consistency and real-time feasibility to be verified by the same C kernel.}



\section{Experiments and Results}

\CNEN{在所有实验中，模型基于 TensorFlow 和 Adam 优化器进行训练，batch size 设为 \valBatchSize，总训练轮次为 \valEpochs。模型采用 \valLearningRateSchedule{}，学习率 $\lr(\lridx)$ 设为 \valInitLR，损失函数中 $\wMAE$ 取值为 \valMaeWeight。}{In all experiments, models were trained with TensorFlow and the Adam optimizer, with batch size set to \valBatchSize{} and total training epochs set to \valEpochs. The model adopted \valLearningRateSchedule{}, the learning rate $\lr(\lridx)$ was set to \valInitLR, and $\wMAE$ in the loss function was set to \valMaeWeight.%<CNHash>2F826926</CNHash>
}

\CNEN{本文采用三项物理指标评价补偿质量，并通过板端数值一致性、实测吞吐率和内存占用评价部署可行性。频率漂移定义为 10--128~Hz 拟合固有频率在不同震级扫描条件下相对于中位数的最大变化量。灵敏度漂移定义为 \valSensitivityCompareFreq~Hz 处插值灵敏度在不同震级条件下相对于中位数的最大变化量，单位为 $\text{V}\cdot\text{s}/\text{m}$。线性度误差定义为 128~Hz 以下输入--输出幅值拟合中 $1-R^2$ 的平均百分比。KEIL speed 是目标板实测吞吐率，QEMU-MAE 和 KEIL-MAE 用于验证导出的 C 实现与 TensorFlow 参考输出之间的数值一致性。}{Three physical metrics are used to evaluate compensation quality, while deployment feasibility is assessed via on-board numerical consistency, measured throughput, and memory footprint. Frequency drift is defined as the maximum deviation of the fitted natural frequency from its median across different excitation levels within 10--128~Hz. Sensitivity drift is defined as the maximum deviation of the interpolated sensitivity at \valSensitivityCompareFreq~Hz from its median across different excitation levels, with units of $\text{V}\cdot\text{s}/\text{m}$. Linearity error is defined as the mean percentage of $1-R^2$ in the input-output amplitude fitting below 128~Hz. KEIL speed is the measured on-board throughput, while QEMU-MAE and KEIL-MAE verify numerical consistency between the exported C implementation and the TensorFlow reference output.%<CNHash>8D609968</CNHash>
}

\subsection{Wiener-KAN Compensation}

\CNEN{本节集中评价所提出的 Wiener-KAN 补偿模型。未补偿 MET 响应在同一频率--震级矩阵上表现出 \valMainWienerFreqDrift~Hz 频率漂移、\valMainWienerSensDrift~$\text{V}\cdot\text{s}/\text{m}$ 灵敏度漂移和 \valMainWienerLinearity\% 线性度误差；Wiener-KAN 将三项指标分别降低到 \valPrimaryFreqDrift~Hz、\valPrimarySensDrift~$\text{V}\cdot\text{s}/\text{m}$ 和 \valPrimaryLinearity\%。}{This section focuses on evaluating the proposed Wiener-KAN compensation model. The uncompensated MET response exhibits a frequency drift of \valMainWienerFreqDrift~Hz, a sensitivity drift of \valMainWienerSensDrift~$\text{V}\cdot\text{s}/\text{m}$, and a linearity error of \valMainWienerLinearity\% over the same frequency--magnitude matrix; Wiener-KAN reduces these three metrics to \valPrimaryFreqDrift~Hz, \valPrimarySensDrift~$\text{V}\cdot\text{s}/\text{m}$, and \valPrimaryLinearity\%, respectively.%<CNHash>E6CDA30F</CNHash>
}

\CNEN{三项指标覆盖了 MET 大震级非线性的主要外在表现。频率漂移反映共振位置随输入震级迁移的程度，灵敏度漂移反映固定频点处通带增益的幅值依赖性，线性度误差反映输入--输出幅值关系在带内频率范围内偏离线性参考的程度。Wiener-KAN 在三项指标上的同步下降说明，补偿结果在多个频点和多个波形片段上保持一致改善，并在频率--震级矩阵上整体收束。}{The three metrics cover the main observable manifestations of large-magnitude MET nonlinearity. Frequency drift reflects the migration of the resonance position with input magnitude, sensitivity drift reflects the magnitude dependence of passband gain at a fixed frequency, and linearity error reflects the deviation of the input-output amplitude relation from the linear reference within the in-band range. The simultaneous reduction of all three metrics indicates that Wiener-KAN improves the response across multiple frequency points and waveform segments over the frequency--magnitude matrix.}

\begin{figure*}[!t]
    \centering
    \includegraphics[width=0.85\textwidth]{../../../ex_projects/plot/multi/fig_08_frequency_response_comparison/data/fig_08_frequency_response_comparison.png}
    \caption{Effectiveness of the proposed Wiener-KAN compensation model in reducing amplitude-dependent frequency-response drift. (a) Frequency-response curves before and after compensation under representative magnitude conditions. (b) Natural-frequency trajectory versus magnitude before and after compensation. (c) Sensitivity trajectory at \valSensitivityCompareFreq~Hz before and after compensation. (d) Representative time-domain outputs under selected frequency and magnitude conditions. (e) Frequency-drift and sensitivity-drift values before and after compensation.}
    \label{fig:frequency_response_comparison}
\end{figure*}

\CNEN{Fig.~\ref{fig:frequency_response_comparison}(a) 展示补偿前后的幅频响应曲线，显示不同震级条件下的响应包络明显收束。Fig.~\ref{fig:frequency_response_comparison}(b)--(c) 展示了 Wiener-KAN 对幅值相关频率响应漂移的抑制效果。补偿后，不同震级下的固有频率轨迹和 \valSensitivityCompareFreq~Hz 灵敏度轨迹向低幅值参考响应收敛。Fig.~\ref{fig:frequency_response_comparison}(d) 给出代表性时域波形，预测波形在相位和幅值上与理想参考输出保持一致。Fig.~\ref{fig:frequency_response_comparison}(e) 给出频率漂移和灵敏度漂移的总体下降幅度。}{Fig.~\ref{fig:frequency_response_comparison}(a) shows the amplitude-frequency response curves before and after compensation, where the response envelopes under different excitation levels are significantly narrowed. Fig.~\ref{fig:frequency_response_comparison}(b)--(c) illustrate the suppression of amplitude-dependent frequency response drift by Wiener-KAN. After compensation, the natural frequency trajectories and sensitivity trajectories at \valSensitivityCompareFreq~Hz converge toward the low-amplitude reference response. Fig.~\ref{fig:frequency_response_comparison}(d) presents representative time-domain waveforms, where predicted waveforms match the ideal reference output in both phase and amplitude. Fig.~\ref{fig:frequency_response_comparison}(e) shows the overall reduction in frequency drift and sensitivity drift.%<CNHash>A4608591</CNHash>
}

\CNEN{从曲线形态看，补偿后的响应保留了低幅值参考系统的主要带通形状，同时压缩了由大震级输入引起的峰值位置和增益差异。时域波形对齐进一步说明，频域收束没有通过牺牲相位一致性获得。该结果与方法部分的设计目标一致：Wiener 线性层负责引入局部频率响应先验，KAN 映射负责校正剩余幅值相关非线性。}{From the curve shape, the compensated response preserves the main bandpass shape of the low-amplitude reference system while reducing the peak-position and gain differences induced by large-magnitude inputs. The time-domain waveform alignment further indicates that the frequency-domain convergence is achieved while preserving phase consistency. This result is consistent with the method design: the Wiener linear layer introduces the local frequency response prior, and the KAN mapping corrects the remaining magnitude-dependent nonlinearity.}

\subsection{Comprehensive Comparison with Baseline Models}

\CNEN{为评估 Wiener-KAN 的综合性能，本文将其与多类具有代表性的非线性补偿方法进行比较，包括传统 Wiener 模型，以及 1DCNN、TCN、RNN、GRU、LSTM 和 LSTMTransformer 等时序学习基线。各基线采用与所提 Wiener-KAN 相同的数据划分、训练轮次和固定学习率策略，并按相近参数规模或资源约束选择宽度、深度、卷积核和循环单元数量，以保持基线比较的规模可比性。}{To evaluate the comprehensive performance of Wiener-KAN, this paper compares it with representative nonlinear compensation methods, including the conventional Wiener model and time-series learning baselines such as 1DCNN, TCN, RNN, GRU, LSTM, and LSTM-Transformer. All baselines adopt the same data split, training epochs, and fixed learning rate strategy as the proposed Wiener-KAN, with width, depth, convolution kernel, and recurrent unit counts selected under comparable parameter budgets or resource constraints to maintain scale comparability in the baseline comparison.%<CNHash>B3308C07</CNHash>
}

\begin{table*}[ht]
    \centering
    \small
    \setlength{\tabcolsep}{2pt}
    \caption{Main comparison under the physical-calibration and on-board efficiency protocol. Lower values are better for frequency drift, sensitivity drift, linearity error, and RAM; higher values are better for KEIL speed.}
    \label{tab:models}
    \begin{tabular}{p{0.16\textwidth}p{0.14\textwidth}p{0.14\textwidth}p{0.14\textwidth}p{0.18\textwidth}p{0.10\textwidth}}
        \hline\hline
        \textbf{Method} & \makecell{\textbf{Freq}\\\textbf{drift}\\(Hz)} & \makecell{\textbf{Sens}\\\textbf{drift}\\($\text{V}\cdot\text{s}/\text{m}$)} & \makecell{\textbf{Linearity}\\(\%)} & \makecell{\textbf{KEIL speed}\\(points/s)} & \makecell{\textbf{RAM}\\(KB)} \\
        \hline
        Wiener & \valMainWienerFreqDrift & \valMainWienerSensDrift & \valMainWienerLinearity & -- & -- \\
        1DCNN & \valMainCNNFreqDrift & \valMainCNNSensDrift & \valMainCNNLinearity & \valMainCNNKeilFps & \valMainCNNKeilRamKB \\
        TCN & \valMainTCNFreqDrift & \valMainTCNSensDrift & \valMainTCNLinearity & \valMainTCNKeilFps & \valMainTCNKeilRamKB \\
        RNN & \valMainRNNFreqDrift & \valMainRNNSensDrift & \valMainRNNLinearity & \valMainRNNKeilFps & \valMainRNNKeilRamKB \\
        GRU & \valMainGRUFreqDrift & \valMainGRUSensDrift & \valMainGRULinearity & \valMainGRUKeilFps & \valMainGRUKeilRamKB \\
        LSTM & \valMainLSTMFreqDrift & \valMainLSTMSensDrift & \valMainLSTMLinearity & \valMainLSTMKeilFps & \valMainLSTMKeilRamKB \\
        LSTM-Trans. & \valMainLSTMTransformerFreqDrift & \valMainLSTMTransformerSensDrift & \valMainLSTMTransformerLinearity & \valMainLSTMTransformerKeilFps & \valMainLSTMTransformerKeilRamKB \\
        \textbf{Wiener-KAN} & \textbf{\valPrimaryFreqDrift} & \textbf{\valPrimarySensDrift} & \textbf{\valPrimaryLinearity} & \textbf{\valPrimaryKeilFps} & \textbf{\valPrimaryKeilRamKB} \\
        \hline\hline
    \end{tabular}
\end{table*}

\CNEN{Tab.~\ref{tab:models} 在统一 MET 补偿协议下比较了 Wiener-KAN 与典型时序基线。Wiener-KAN 同时取得最低频率漂移、最低灵敏度漂移和最高板端吞吐率；TCN 在线性度误差上接近 Wiener-KAN，但其 KEIL speed 低于所提模型；RNN 的板端吞吐率较高，但物理补偿指标的抑制效果受限。}{Tab.~\ref{tab:models} compares Wiener-KAN with typical time-series baselines under a unified MET compensation protocol. Wiener-KAN simultaneously achieves the lowest frequency drift, lowest sensitivity drift, and highest on-board throughput; TCN approaches Wiener-KAN in linearity error, but its KEIL speed is lower than that of the proposed model; RNN offers high on-board throughput, yet its suppression of physical compensation metrics is limited.%<CNHash>D9EDB633</CNHash>
}

\CNEN{该比较表明，物理补偿精度和板端效率需要同时考察。卷积和循环模型能够学习一定的时序相关性，但其频率响应骨架主要由数据驱动优化形成，因此在跨震级频率响应保持方面容易受到训练样本分布影响。Wiener-KAN 通过显式局部传递函数初始化降低了频域学习负担，并通过 LUT 推理减少在线计算量，使物理指标和目标板吞吐率同时保持优势。}{This comparison shows that physical compensation accuracy and on-board efficiency must be examined together. Convolutional and recurrent models can learn temporal correlations, but their frequency response skeletons are mainly formed by data-driven optimization and are therefore sensitive to the training-sample distribution when maintaining frequency response across magnitudes. Wiener-KAN reduces the frequency-domain learning burden through explicit local transfer-function initialization and reduces online computation through LUT inference, maintaining advantages in both physical metrics and target-board throughput.}

\begin{figure*}[ht]
    \centering
    \includegraphics[width=0.98\textwidth]{../../../ex_projects/plot/multi/fig_02_horizontal_summary/data/fig_02_horizontal_summary.png}
    \caption{Main benchmark summary. (a) Natural-frequency range across magnitudes. (b) Sensitivity range at \valSensitivityCompareFreq~Hz across magnitudes. (c) In-band linearity-error range across frequency points. (d) STM32F405 measured KEIL throughput with exported-kernel error annotations. (e) Five-metric radar comparison across physical accuracy and on-board efficiency metrics. (f) Main-benchmark validation-loss trajectories. Error bars in (a)--(c) indicate the corresponding min--max ranges.}
    \label{fig:horizontal_summary}
\end{figure*}

\CNEN{Fig.~\ref{fig:horizontal_summary}(a)--(c) 展示了代表性模型的固有频率、\valSensitivityCompareFreq~Hz 灵敏度和带内线性度误差范围，Wiener-KAN 在固有频率和灵敏度分布上最为收敛。Fig.~\ref{fig:horizontal_summary}(d)--(f) 分别展示板端实测吞吐率、五指标雷达图和基线比较训练收敛曲线。五指标雷达图将频率漂移、灵敏度漂移、线性度误差、KEIL speed 和 KEIL RAM 归一化到共同坐标中比较，显示 Wiener-KAN 在物理补偿和嵌入式效率之间保持更均衡的指标形状。}{Fig.~\ref{fig:horizontal_summary}(a)--(c) show the natural frequency, sensitivity at \valSensitivityCompareFreq~Hz, and in-band linearity error ranges of representative models. Wiener-KAN exhibits the most convergent distributions in natural frequency and sensitivity. Fig.~\ref{fig:horizontal_summary}(d)--(f) present the on-board measured throughput, a five-metric radar chart, and the training convergence curves of the baseline comparison, respectively. The radar chart normalizes frequency drift, sensitivity drift, linearity error, KEIL speed, and KEIL RAM onto common coordinates, showing that Wiener-KAN maintains a more balanced metric profile between physical compensation and embedded efficiency.%<CNHash>E52D85FF</CNHash>
}

\CNEN{训练收敛曲线也支持这一结论。各模型使用一致的数据划分和训练设置，因而 loss 轨迹主要反映结构对同一补偿任务的适配程度。Wiener-KAN 的收敛结果与物理指标保持一致，说明局部频率响应先验、带符号与对称性条件的 KAN 和联合损失共同作用于最终补偿质量。}{The training convergence curves support this conclusion. Since all models use the same data split and training setting, the loss trajectories mainly reflect how each structure adapts to the same compensation task. The convergence result of Wiener-KAN is consistent with the physical metrics, indicating that the local frequency response prior, sign- and symmetry-conditioned KAN, and joint loss act together on the final compensation quality.}

\CNEN{从部署角度看，横向比较还说明模型精度不能脱离导出实现评价。部分时序模型在物理指标上具有一定补偿能力，但其在线路径依赖循环状态更新或卷积计算，目标板吞吐率和内存占用会限制其用于连续采样。Wiener-KAN 的 LUT 路径将主要非线性计算前移到离线表生成阶段，使同一模型同时具备频响补偿能力和低延迟执行能力。}{From the deployment perspective, the horizontal comparison also shows that model accuracy must be evaluated together with the exported implementation. Some time-series models provide a degree of compensation in the physical metrics, but their online paths depend on recurrent state updates or convolutional computation, so target-board throughput and memory footprint constrain their use in continuous sampling. The LUT path of Wiener-KAN moves the main nonlinear computation to offline table generation, allowing the same model to combine frequency-response compensation with low-latency execution.}

\subsection{Ablation Studies of Wiener-KAN Components}

\CNEN{消融实验评价 Wiener-KAN 的训练目标、结构变体、符号与对称性条件、Wiener 前端和超参数设置。Tab.~\ref{tab:ablation_overview} 汇总 \valAblationOverviewRowCount{} 条数值对比，覆盖结构变体、符号与对称性条件、Wiener 前端和损失函数。}{The ablation study evaluates the training objective, structural variants, sign and symmetry conditions, Wiener front-end, and hyperparameter settings of Wiener-KAN. Tab.~\ref{tab:ablation_overview} summarizes \valAblationOverviewRowCount{} numerical comparisons covering structural variants, sign and symmetry conditions, the Wiener front-end, and the loss function.%<CNHash>9B63CC30</CNHash>
}

\begin{table*}[ht]
    \centering
    \small
    \setlength{\tabcolsep}{2pt}
    \caption{Numerical overview of Wiener-KAN ablation studies. Each row reports the evaluated physical-metric triplet for one design contrast.}
    \label{tab:ablation_overview}
    \resizebox{\textwidth}{!}{% Auto-generated table; do not edit manually.
\begingroup
\renewcommand{\arraystretch}{1.22}
\begin{tabular}{@{}p{0.12\textwidth}p{0.18\textwidth}p{0.40\textwidth}ccc@{}}
\hline\hline
\textbf{Group} & \textbf{Variant} & \textbf{Controlled contrast} & \makecell{\textbf{Freq}\\\textbf{drift}\\(Hz)} & \makecell{\textbf{Sens}\\\textbf{drift}\\(\%)} & \makecell{\textbf{Linearity}\\(\%)} \\ \hline
\textbf{Baseline} & \textbf{B-spline + Odd+pos. + System, fixed, + MAE + AFMAE} & \textbf{Shared canonical Wiener-KAN setting} & \textbf{2.34} & \textbf{9.09} & \textbf{0.511} \\
\hdashline
\multirow{2}{=}{Struct} & CNN-KAN & Conv front replaces Wiener front & 2.43 & 9.22 & 0.883 \\
 & Wiener-MLP & MLP replaces KAN mapping & 5.82 & 27.48 & 0.782 \\
\hdashline
\multirow{2}{=}{Constr.} & Positive & No odd symmetry & 10.91 & 10.61 & 2.109 \\
 & Odd & No positive basis & 58.03 & 50.87 & 38.586 \\
\hdashline
\multirow{3}{=}{Wiener front} & Random / fixed & Random init., frozen & 5.37 & 22.28 & 1.308 \\
 & System / trainable & Measured init., trainable & 94.87 & 64.68 & 7.478 \\
 & Random / trainable & Random init., trainable & 102.57 & 16.65 & 5.144 \\
\hdashline
\multirow{2}{=}{Loss} & MAE & Time-domain pointwise loss & 13.00 & 20.28 & 0.793 \\
 & AFMAE & Amplitude-frequency loss & 4.00 & 21.06 & 1.068 \\
\hline\hline
\end{tabular}
\endgroup
}
\end{table*}

\CNEN{结构消融表明，CNN-KAN 和 Wiener-MLP 分别对应 Wiener-KAN 的线性动态骨架和非线性补偿主体。CNN-KAN 在频率漂移与灵敏度漂移上接近基线配置，但在线性度误差上表现出明显的误差上升，说明局部频率响应初始化的 Wiener 前端更有利于保持带内补偿曲线的稳定形状。Wiener-MLP 在三项物理指标上均落后于 Wiener-KAN，说明共享 Wiener 前端后，带物理约束的 KAN 映射仍然承担幅值相关非线性补偿的主要表达能力。}{The structural ablation indicates that CNN-KAN and Wiener-MLP correspond to the linear dynamic skeleton and nonlinear compensation body of Wiener-KAN, respectively. CNN-KAN approaches the baseline configuration in frequency drift and sensitivity drift, yet exhibits a marked increase in linearity error, suggesting that the Wiener front-end initialized with local frequency response is more conducive to preserving the stable shape of the in-band compensation curve. Wiener-MLP lags behind Wiener-KAN on all three physical metrics, indicating that after sharing the Wiener front-end, the physics-constrained KAN mapping still undertakes the primary expressive capability for amplitude-dependent nonlinear compensation.%<CNHash>B3ED615F</CNHash>
}

\CNEN{约束消融实验将奇对称性与正值性分开考察。Tab.~\ref{tab:ablation_overview} 中的 \textit{Odd / No positive basis} 对照保留奇对称性，但移除了正输入区间的方向保持约束，其频率漂移、灵敏度漂移和线性度误差同时上升，说明正值约束为 Wiener-KAN 提供了稳定的函数形态先验。Wiener front 分组进一步表明，冻结的系统先验设置给出最好的频率漂移、灵敏度漂移和线性度误差，而可训练 IIR 前端导致频率漂移指标明显上升。}{The constraint ablation study examines odd symmetry and positivity separately. The \textit{Odd / No positive basis} variant in Tab.~\ref{tab:ablation_overview} retains odd symmetry but removes the direction-preserving constraint on the positive input interval, increasing frequency drift, sensitivity drift, and linearity error simultaneously; this indicates that the positivity constraint provides Wiener-KAN with a stable functional shape prior. The Wiener front group further shows that the frozen system prior setting yields the best frequency drift, sensitivity drift, and linearity error, while the trainable IIR front-end causes a marked increase in frequency drift.%<CNHash>4C89F1CB</CNHash>
}

\CNEN{损失函数消融显示，MAE+AFMAE 在频率漂移、灵敏度漂移和线性度误差三项指标上形成最均衡结果，说明时域逐点误差和幅频响应误差需要共同参与训练目标。单独使用 MAE 时，逐点误差目标对幅值相关频率漂移的约束较弱；单独使用 AFMAE 时，幅频统计约束改善频率漂移，但灵敏度漂移和线性度误差仍高于联合目标。结构变体、物理约束、Wiener 前端、超参数敏感性和损失函数消融的详细结果见 Supplementary Note 7。}{The loss-function ablation shows that MAE+AFMAE yields the most balanced results across frequency drift, sensitivity drift, and linearity error, indicating that pointwise time-domain error and amplitude-frequency response error must jointly participate in the training objective. With MAE alone, the pointwise error objective imposes weak constraints on amplitude-related frequency drift. With AFMAE alone, the amplitude-frequency statistical constraint improves frequency drift, but sensitivity drift and linearity error remain higher than those of the joint objective. Detailed structural, constraint, Wiener-front-end, hyperparameter, and loss-function ablations are provided in Supplementary Note 7.%<CNHash>A8825545</CNHash>
}

\CNEN{这些消融结果共同限定了各组件的作用范围。Wiener 前端主要决定频域动态骨架，KAN 映射主要决定幅值相关非线性表达，符号与对称性条件稳定函数形态，MAE+AFMAE 则把时域波形误差和幅频误差同时纳入训练。任一组件被替换或移除时，至少一项物理指标出现上升，说明结构先验、函数条件和训练目标的组合共同决定最终性能。}{These ablation results jointly define the role of each component. The Wiener front-end mainly determines the frequency-domain dynamic skeleton, the KAN mapping mainly determines the magnitude-dependent nonlinear expression, the sign and symmetry conditions stabilize the function shape, and MAE+AFMAE brings both time-domain waveform error and amplitude-frequency error into training. When any component is replaced or removed, at least one physical metric increases, showing that the final performance is jointly determined by the structural prior, functional conditions, and training objective.}

\CNEN{板端推理测试确认 Wiener-KAN 从训练模型到嵌入式 C 实现的完整链路。训练完成后，权重、归一化常数和 LUT 表被导出为 C 代码；同一份 C 实现先在 QEMU 中运行并与 TensorFlow 参考输出比较，再烧录到 \valTargetMcu{} 目标板，通过串口基准循环采集吞吐率，并从 Keil 构建结果中读取 Flash 与 RAM 占用。Wiener-KAN 主实现达到 \valPrimaryKeilFps{} points/s，RAM 占用为 \valPrimaryKeilRamKB~KB。详细验证流程、LUT 变体和量化点 sweep 见 Supplementary Note 6。}{On-board inference testing verifies the complete pipeline of Wiener-KAN from the trained model to the embedded C implementation. After training, the weights, normalization constants, and LUT tables are exported as C code. The same C implementation is first run in QEMU and compared with the TensorFlow reference output, then flashed onto the \valTargetMcu{} target board. Throughput is collected via a serial benchmark loop, and Flash and RAM usage are read from the Keil build results. The primary Wiener-KAN implementation achieves \valPrimaryKeilFps{} points/s with \valPrimaryKeilRamKB~KB RAM. The detailed verification flow, LUT variants, and quantization-point sweep are provided in Supplementary Note 6.%<CNHash>CD1E9BC4</CNHash>
}
\FloatBarrier

\CNEN{板端结果为实时前馈补偿提供了直接证据。QEMU 与 TensorFlow 的一致性检查用于排除导出内核的数值偏差，目标板串口基准则测量包含存储访问和函数调用在内的实际运行速度。由于 LUT 路径把样条函数求值转化为表访问，主实现能够在较小 RAM 占用下达到高吞吐率，从而满足连续采样链路中的逐点补偿目标。}{The on-board results provide direct evidence for real-time feedforward compensation. The QEMU-to-TensorFlow consistency check excludes numerical deviation in the exported kernel, while the target-board serial benchmark measures the actual runtime speed including memory access and function-call overhead. Because the LUT path converts spline-function evaluation into table access, the primary implementation achieves high throughput with a small RAM footprint, satisfying the pointwise compensation target in a continuous sampling chain.}



\section{Discussion}

\CNEN{本文验证基于一个 MET 样品、受控正弦激励矩阵以及 $\valEnvironmentTemperature$ 环境。该设置提供了可复现的频率--震级响应矩阵，用于量化震级相关漂移；更多温度条件、更多传感器样品以及接近野外地震记录的复合波形需要独立测量数据支撑。嵌入式结果通过 QEMU 数值检查和 \valTargetMcu{} 板端测量验证导出实现。KEIL speed 和 KEIL RAM Usage 描述当前编译器、MCU、数值精度和窗口长度下的实现结果，QEMU-MAE 与 KEIL-MAE 用于检查导出 C 实现与 TensorFlow 参考输出之间的数值一致性。其他 MCU 系列、定点实现、多通道采集和不同编译设置会改变实测吞吐率与内存占用。}{The validation in this paper is based on one MET sensor sample, a controlled sinusoidal excitation matrix, and an ambient temperature of $\valEnvironmentTemperature$. This setup provides a reproducible frequency--magnitude response matrix for quantifying magnitude-dependent drift; additional temperature conditions, more sensor samples, and composite waveforms resembling field seismic records require independent measurement data. The embedded results are verified through QEMU numerical checks and \valTargetMcu{} on-board measurements of the exported implementation. KEIL speed and KEIL RAM Usage describe the implementation results under the current compiler, MCU, numerical precision, and window length; QEMU-MAE and KEIL-MAE are used to check numerical consistency between the exported C implementation and the TensorFlow reference output. Other MCU families, fixed-point implementations, multi-channel acquisition, and different compilation settings will alter the measured throughput and memory footprint.%<CNHash>249D6DD6</CNHash>
}

\CNEN{这些适用条件保留本文关于方法可行性的核心结论，并限定了结果的适用范围。本文给出的频率漂移、灵敏度漂移和线性度误差应理解为当前样品、当前频率--震级矩阵和当前目标参考系统下的补偿结果；板端速度和内存占用应理解为当前 \valTargetMcu{}、浮点实现和导出内核下的实现结果。将该流程迁移到其他 MET 结构时，需要重新测量局部频率响应、重建目标参考系统，并按相同指标重新评价。}{These applicability conditions preserve the core conclusion on method feasibility and delimit the scope of the results. The reported frequency drift, sensitivity drift, and linearity error should be interpreted as compensation results under the current sample, frequency--magnitude matrix, and target reference system; the on-board speed and memory footprint should be interpreted as implementation results under the current \valTargetMcu{}, floating-point realization, and exported kernel. When transferring the workflow to other MET structures, the local frequency responses should be remeasured, the target reference system rebuilt, and the same metrics recomputed.}

\subsection{Data Availability}

\CNEN{现有公开传感器漂移数据集主要覆盖其他传感器类型，例如气体传感器阵列\cite{vergaraChemicalGasSensor2012}。据我们所知，公开资源尚未提供面向幅值相关时域补偿的同步 MET 输出、参考激励、频率、震级和目标标定曲线。本文使用的数据集是面向 MET 漂移补偿自建的频率--震级响应矩阵。处理后的响应矩阵、代表性波形、数据划分脚本和指标计算脚本公开于 \texttt{github.com/pikastech/wiener-kan}.}{Existing public sensor drift datasets mainly cover other sensor types, such as gas sensor arrays \cite{vergaraChemicalGasSensor2012}. To our knowledge, no public resource provides synchronized MET output, reference excitation, frequency, magnitude, and target calibration curves for amplitude-dependent time-domain compensation. The dataset used in this paper is a self-built frequency--magnitude response matrix for MET drift compensation. The processed response matrix, representative waveforms, data partitioning scripts, and metric calculation scripts are available at \texttt{github.com/pikastech/wiener-kan}.%<CNHash>133F80DB</CNHash>
}
\FloatBarrier

\section{Conclusion}

\CNEN{本文结合 Wiener 结构的线性--非线性分解、频率响应信息和 Kolmogorov-Arnold 网络，构建了 Wiener-KAN 非线性补偿系统。本文构建并开放了捕获幅频耦合效应的三维频率响应数据集，将相对频率响应切片映射到 Wiener 线性层初始化权重，使用带符号和对称性约束的可训练样条激活函数直接学习传感器校准曲线，并通过 AFMAE 约束幅频差异。实验结果表明，Wiener-KAN 改善了电化学地震检波器的输入输出曲线并抑制频率漂移，说明前馈补偿可以在保持传感器硬件结构的同时改善 MET 传感器的大震级响应质量。基于预计算查找表的推理机制进一步验证了地震仪器板端实时计算的可行性。}{This work constructs the Wiener-KAN nonlinear compensation system by combining the linear-nonlinear decomposition of the Wiener structure, frequency response information, and Kolmogorov-Arnold networks. This paper constructs and releases a three-dimensional frequency response dataset capturing amplitude-frequency coupling effects, maps relative frequency response slices into the initialization weights of the Wiener linear layer, uses trainable spline activation functions with sign and symmetry constraints to directly learn the sensor calibration curve, and constrains amplitude-frequency discrepancy through AFMAE. Experimental results show that Wiener-KAN improves the input-output curve of electrochemical seismometers and suppresses frequency drift, indicating that feedforward compensation can improve the large-magnitude response quality of MET sensors while preserving the sensor hardware structure. The inference mechanism based on a precomputed lookup table further demonstrates the feasibility of real-time on-board computation in seismic instruments.%<CNHash>55797861</CNHash>
}

\CNEN{在完整流程层面，本文将幅频耦合测试、局部频率响应先验、带符号和对称性条件的 KAN 非线性映射、AFMAE 训练目标以及基于 LUT 的板端推理连接为闭合工作流。该工作流保留了 MET 小信号频率响应模型中可解释的部分，并将解析困难的震级相关非线性分配给可训练函数。Wiener-KAN 因此为电化学地震检波器在不增加力反馈硬件的情况下实现大震级实时补偿提供了可复现路径。}{At the full-pipeline level, this paper connects amplitude-frequency coupling tests, local frequency response priors, sign- and symmetry-conditioned KAN nonlinear mapping, the AFMAE training objective, and LUT-based on-board inference into a closed workflow. The workflow preserves the interpretable part of the MET small-signal frequency response model while assigning the analytically difficult magnitude-dependent nonlinearity to trainable functions. Wiener-KAN therefore provides a reproducible path for real-time large-magnitude compensation of electrochemical seismometers without adding force-feedback hardware.}

\FloatBarrier

\section*{\CNEN{补充材料}{Supplementary Information}}
\CNEN{补充材料随稿件发布于 Microsystems \& Nanoengineering 网站（http://www.nature.com/micronano）。}{Supplementary information accompanies the manuscript on the Microsystems \& Nanoengineering website (http://www.nature.com/micronano).}

\section*{Funding}
This work was funded by the Deep Earth Probe and Mineral Resources Exploration-National Science and Technology Major Project (Grant No. 2024ZD1002503).

\FloatBarrier

\bibliography{nonlinear}

\end{document}


